\documentclass[11pt,a4paper]{report}
\usepackage[T1]{fontenc}
\usepackage{polski}
\usepackage[utf8]{inputenc}
\usepackage[margin=3.5cm]{geometry}
\usepackage{graphicx}
\usepackage[hidelinks]{hyperref}

\setcounter{tocdepth}{1}
\setcounter{secnumdepth}{1}

\begin{document}

\begin{titlepage}
\centering
\IfFileExists{logo_AGH.pdf}{%
  \includegraphics[width=\textwidth,height=3cm,keepaspectratio]{logo_AGH.pdf}
}{%
  {\Large\sffamily Akademia Górniczo-Hutnicza w Krakowie}
}
\par\vspace{5mm}
{\bfseries\sffamily WYDZIAŁ ODLEWNICTWA\par}
\vspace{5mm}
{\bfseries\sffamily KATEDRA ---\par} % Uzupełnij nazwę katedry.
\vspace{14mm}
{\sffamily\huge Projekt dyplomowy\par}
\vspace{12mm}
{\sffamily\Large
Projekt i implementacja systemu cyfrowego bliźniaka wózka AGV
z wykorzystaniem architektury Edge Computing\par}
\vspace{5mm}
{\sffamily\large
Design and implementation of a digital twin system for an AGV trolley
using the Edge Computing architecture\par}
\vfill
\begin{flushleft}
\sffamily
\begin{tabular}{@{}ll@{}}
Autor: & Patryk Pięta\\
Kierunek studiów: & Technologie przemysłu 4.0\\
Opiekun pracy: & dr inż. Waldemar Bauer\\
\end{tabular}
\end{flushleft}
\vspace{10mm}
{\sffamily Kraków, 2026\par}
\end{titlepage}

% Robocowy szkielet pracy i propozycja spisu treści.
% Przed oddaniem dodaj wymagane przez wydział oświadczenia,
% streszczenie po polsku i po angielsku oraz opis wykorzystania AI,
% zgodnie z aktualnym wzorem otrzymanym od uczelni.
% Nie zakładaj, że sam obraz AI_poziom_3.jpg stanowi oświadczenie.

\pagenumbering{roman}
\tableofcontents
\clearpage
\pagenumbering{arabic}

\chapter{Wstęp}
\section{Uzasadnienie wyboru tematu}
% Potrzeba obserwacji stanu AGV i powiązanie z Przemysłem 4.0.
\section{Cel i zakres pracy}
% Cel: zaprojektowanie, implementacja i weryfikacja prototypu
% odwzorowującego stan AGV na podstawie strumienia telemetrii.
% Określ granice prototypu i mierzalne kryteria realizacji celu.
\section{Wkład własny autora}
% Wskaż samodzielnie wykonane elementy projektu oraz źródła
% wykorzystanych modeli, bibliotek i innych materiałów.
\section{Struktura pracy}

\chapter{Podstawy teoretyczne i przegląd literatury}
\section{Wózki AGV i monitorowanie ich stanu}
\section{Cyfrowy model, cyfrowy cień i cyfrowy bliźniak}
% Wyjaśnij kierunek przepływu danych i określ, w jakim zakresie
% prototyp realizuje koncepcję cyfrowego bliźniaka.
\section{Architektura Edge Computing w systemach przemysłowych}
\section{Przetwarzanie strumieniowe i przechowywanie szeregów czasowych}
\section{Przegląd podobnych rozwiązań i uzasadnienie przyjętego podejścia}
% Porównuj rozwiązania według kryteriów związanych z celem pracy.

\chapter{Założenia i wymagania projektowe}
\section{Przeznaczenie systemu i scenariusze użytkowania}
\section{Wymagania funkcjonalne}
% Telemetria, zapis historii, wizualizacja, LiDAR i ostrzeżenia.
\section{Wymagania niefunkcjonalne i kryteria akceptacji}
% Określ sprawdzalne wymagania dotyczące aktualności danych,
% niezawodności, wydajności i możliwości uruchomienia systemu.
\section{Zakres danych oraz założenia dotyczące modelu AGV}
\section{Ograniczenia sprzętowe i programowe}
% Rozróżnij dane syntetyczne i dane z rzeczywistych czujników.
% Obecna wersja nie steruje fizycznym wózkiem.

\chapter{Projekt architektury systemu}
\section{Podział na warstwę Edge i usługi centralne}
% Schemat komponentów i rzeczywiste miejsce wykonywania obliczeń.
\section{Dobór technologii i porównanie alternatyw}
% Uzasadnij użycie Kafka, TimescaleDB, Spark, Grafana i Three.js.
\section{Przepływ danych i interfejsy komunikacyjne}
\section{Schemat telemetrii i model przechowywania danych}
\section{Model ruchu i aktualizacja wizualizacji AGV}
% Opisz estymację ruchu w przeglądarce i jej ograniczenia.
\section{Obsługa błędów, trwałość danych i założenia bezpieczeństwa}

\chapter{Implementacja systemu}
\section{Emulator telemetrii i scenariusze pracy wózka}
\section{Przetwarzanie danych LiDAR w warstwie Edge}
% Filtracja i grupowanie DBSCAN; obecnie źródłem jest MockRPLidar.
\section{Przesyłanie i walidacja danych w Apache Kafka}
\section{Zapis i agregacja danych w TimescaleDB}
\section{Przetwarzanie strumieniowe i ostrzeżenia w Apache Spark}
% Próg temperatury jest regułą ostrzegania, a nie predykcją awarii.
\section{Wizualizacja 3D i komunikacja przez WebSocket}
\section{Prezentacja danych historycznych w Grafanie}
\section{Konteneryzacja, konfiguracja i panel uruchomieniowy}
% Opisz istotne decyzje i trudności, zamiast omawiać każdy plik.

\chapter{Testy i ocena rozwiązania}
\section{Środowisko testowe i metodyka badań}
% Podaj konfigurację, wersje, czas prób i sposób pomiaru.
\section{Testy poprawności telemetrii, LiDAR i modelu ruchu}
\section{Testy integracyjne przepływu i zapisu danych}
\section{Testy ostrzeżeń i aktualizacji wizualizacji}
\section{Testy obsługi błędów i przerw w dostarczaniu danych}
\section{Pomiary opóźnień i wydajności systemu}
% PLANOWANE: wykonaj pomiary przed opisaniem wyników.
% Odróżnij zaległość konsumenta Kafka od opóźnienia w milisekundach.
% Rozważ częstotliwość odbioru, kompletność danych, CPU i RAM.
\section{Ocena realizacji wymagań i omówienie ograniczeń}
% Powiąż wymagania z wynikami konkretnych testów.
% Nie przedstawiaj testów Raspberry Pi, rzeczywistych czujników
% ani badań długotrwałych jako wykonanych bez ich przeprowadzenia.

\chapter{Podsumowanie i wnioski}
\section{Ocena realizacji celu pracy}
\section{Najważniejsze osiągnięcia i wnioski}
\section{Możliwości zastosowania i kierunki dalszego rozwoju}
% Rzeczywiste czujniki, weryfikacja na Raspberry Pi,
% wspólny stan pojazdu, buforowanie Edge i rozbudowa diagnostyki.

\chapter*{Bibliografia}
\phantomsection
\addcontentsline{toc}{chapter}{Bibliografia}
% Wstaw rzeczywiście wykorzystane publikacje i dokumentację.
% Przy pisaniu pracy zastąp ten nagłówek bibliografią generowaną
% przez BibTeX lub biblatex; nie dodawaj wówczas drugiego nagłówka.

\appendix
\chapter{Instrukcja uruchomienia i obsługi systemu}
\chapter{Specyfikacja komunikatów i konfiguracja usług}
\chapter{Szczegółowe wyniki testów}

\end{document}
